\section{Implementation Environment}

The system was developed and tested using the following technology stack:

\begin{table}[htbp]
\centering
\begin{tabular}{l l}
\toprule
\textbf{Component} & \textbf{Technology} \\
\midrule
Stream Broker       & Apache Kafka 7.5 (via Docker Compose) \\
Coordination        & Apache Zookeeper (via Docker Compose) \\
State Store         & Redis 7 \\
API Server          & FastAPI + Uvicorn (Python 3.11+) \\
Python Worker       & kafka-python-ng, mmh3, psutil, redis-py \\
C++ Worker          & C++20, librdkafka, xxHash, CMake + Ninja \\
Dashboard           & HTML5, CSS3, Chart.js, vanilla JavaScript \\
Containerization    & Docker Compose (multi-service orchestration) \\
Testing             & pytest (28 Python tests), Catch2 (C++ tests) \\
\bottomrule
\end{tabular}
\caption{Technology stack used in the Phase~1 implementation.}
\label{tab:tech_stack}
\end{table}

All infrastructure services (Kafka, Zookeeper, Redis) are orchestrated via a single \texttt{docker-compose.yml} file. The C++ worker can optionally be built and run via Docker using a multi-stage \texttt{Dockerfile.worker}.

\section{Interactive Dashboard}

A fully functional, single-page real-time web dashboard was developed to visualize cluster metrics. The dashboard polls the \texttt{/metrics/cluster} endpoint every 2 seconds and renders all data using Chart.js for interactive charting.

\begin{figure}[htbp]
    \centering
    % PLACEHOLDER: Replace with your dashboard screenshot in Overleaf
    % \includegraphics[width=0.9\textwidth]{Figures/dashboard_screenshot.png}
    \framebox[\textwidth]{\rule{0pt}{7cm}\parbox{0.8\textwidth}{\centering\Large\textbf{[PLACEHOLDER]}\\[1em]\normalsize Upload your dashboard screenshot to the \texttt{Figures/} folder in Overleaf and uncomment the \texttt{includegraphics} line above.}}
    \caption{Web dashboard showing live sketch estimates, heavy hitters list, P99 latency charts, and system resource usage.}
    \label{fig:dashboard_screenshot}
\end{figure}

The dashboard interface is organized into the following sections:

\subsection{Live Sketch Estimate Cards}
Three prominent cards display the live answers from each probabilistic sketch:
\begin{itemize}
    \item \textbf{Unique Items} (HyperLogLog): Estimated number of distinct items with relative error percentage.
    \item \textbf{Top Heavy Hitter} (Misra--Gries): The most frequent item in the current sliding window and its share.
    \item \textbf{Top Item Frequency} (Count-Min Sketch): Estimated frequency count with error bound.
\end{itemize}

\subsection{Key Performance Indicator Strip}
Six compact metric boxes display HLL cardinality, CMS total events processed, P99 latency, average CPU usage, average memory usage, and real-time throughput (events per second).

\subsection{Charts}
\begin{itemize}
    \item \textbf{P99 Latency History:} A multi-line time-series chart showing the latency trend per worker over time. Each worker is rendered in a different color with a pulsing ``LIVE'' indicator.
    \item \textbf{CPU and Memory per Worker:} A grouped bar chart showing both CPU and memory usage side by side for each worker.
\end{itemize}

\subsection{Heavy Hitters Table}
A ranked table of the top-12 most frequent items detected by Misra--Gries, with visual percentage share bars for each entry.

\subsection{Query Explorer}
An interactive panel allowing users to test queries directly from the dashboard:
\begin{itemize}
    \item \textbf{Frequency tab:} Enter an item key and retrieve the CMS frequency estimate with error bounds.
    \item \textbf{Cardinality tab:} Retrieve the current HLL distinct count estimate.
    \item \textbf{Heavy Hitters tab:} Retrieve the top-$N$ items from Misra--Gries.
\end{itemize}

\section{Testing and Validation}

\subsection{Unit Test Suite}
A comprehensive test suite of 28 Python unit tests validates the correctness of all implemented components:

\begin{table}[htbp]
\centering
\begin{tabular}{l c l}
\toprule
\textbf{Component} & \textbf{Tests} & \textbf{Key Assertions} \\
\midrule
Count-Min Sketch & 6 & Error bounds, no-undercount guarantee, memory \\
HyperLogLog      & 7 & Cardinality accuracy, duplicate handling, precision levels \\
Misra--Gries     & 7 & Heavy hitter detection, window eviction, sorting \\
Router           & 8 & Scoring formula, node ranking, edge cases \\
\bottomrule
\end{tabular}
\caption{Summary of the Python unit test suite (28 tests total, all passing).}
\label{tab:test_suite}
\end{table}

All 28 tests pass without requiring any infrastructure (no Kafka or Redis needed). The C++ sketches are independently validated via Catch2 unit tests.

\subsection{Standalone Simulation}
A standalone benchmark (\texttt{scripts/run\_simulation.py}) runs entirely in-process without any external dependencies. It evaluates four experimental configurations across 100,000 synthetic events:

\begin{table}[htbp]
\centering
\begin{tabular}{c l l l}
\toprule
\textbf{Config} & \textbf{Router} & \textbf{Parameters} & \textbf{Description} \\
\midrule
1 & Round-robin  & Static   & Baseline \\
2 & Round-robin  & Adaptive & Controller tunes parameters \\
3 & Smart router & Static   & Load-aware routing only \\
4 & Smart router & Adaptive & Fully adaptive (Phase~2 target) \\
\bottomrule
\end{tabular}
\caption{Four experimental configurations evaluated in the standalone simulation.}
\label{tab:simulation_configs}
\end{table}

\subsection{Sample Query Validation}
The following sample queries were executed against a running system to verify end-to-end correctness:

\begin{lstlisting}[language=bash, caption=Sample frequency query and response.]
$ curl "http://localhost:8000/query/frequency?key=item_1"

{
  "key": "item_1",
  "routed_to": "worker-0",
  "cms_error_bound": 323.31,
  "cms_total_count": 243590,
  "worker_metrics": {
    "template": "medium",
    "memory_percent": 81.7
  }
}
\end{lstlisting}

With a Zipfian distribution ($s = 1.2$), \texttt{item\_1} consistently exhibits the highest estimated frequency, confirming the correctness of the Count-Min Sketch implementation. Similarly, HLL cardinality estimates for a 10,000-item vocabulary consistently fall within the expected $\pm$1.6\% error range for $p = 12$.
